\documentclass[UTF8,10pt,a4paper]{ctexart}
\usepackage[margin=20mm,headheight=26pt]{geometry}
\usepackage{graphicx,xcolor,fancyhdr,booktabs,tabularx,enumitem,hyperref}
\definecolor{ink}{HTML}{12313B}\definecolor{teal}{HTML}{008F89}
\hypersetup{colorlinks=true,urlcolor=teal,linkcolor=teal,pdftitle={WorldGuard 提交材料}}
\setlength{\parindent}{0pt}\setlength{\parskip}{4pt}\linespread{1.12}
\setlist{nosep,leftmargin=1.5em}
\pagestyle{fancy}\fancyhf{}
\fancyhead[L]{\raisebox{-4pt}{\includegraphics[height=8mm]{assets/logo.png}}\quad\textbf{\textcolor{ink}{WorldGuard}}}
\fancyhead[R]{\small 零界智行\quad U406}
\fancyfoot[L]{\footnotesize 本科生组 · 南京大学 / 东南大学}
\fancyfoot[R]{\footnotesize\thepage}
\renewcommand{\headrulewidth}{0.4pt}
\ctexset{section={format=\Large\bfseries\color{ink},beforeskip=12pt,afterskip=6pt},subsection={format=\large\bfseries\color{teal}}}
\newcommand{\cover}[2]{\vspace*{5mm}{\small\textcolor{teal}{WORLDGUARD / SUBMISSION 2026}}\par\vspace{4mm}{\huge\bfseries\color{ink}#1}\par\vspace{3mm}{\large #2}\par\vspace{4mm}\hrule\vspace{5mm}}
\newcommand{\note}[1]{\par\smallskip\noindent\colorbox{teal!7}{\parbox{0.95\linewidth}{\small #1}}\par\smallskip}
\begin{document}
\cover{运行证据核查指南}{真实记录 · 模型参与 · 工具调用 · 独立验证}
\note{阅读顺序：先看下方调用链，再按事件序号查看原记录，最后打开附件 05 复核地图、路径与验收结果。}
\section{AGH 执行记录}
两次真实 Agnes 运行均发生于 2026-10-15。复杂仓库运行包含 3 个案例、19 次模型请求尝试；窄通道运行包含 5 个案例、23 次请求尝试。两个运行子目录分别保存 \texttt{agh\_session.jsonl}、\texttt{agh\_summary.json}、业务事件、决策与结构化报告。

记录为从官方 AGH Host 实际事件捕获的脱敏投影，保留序号、时间、工具 ID、输入输出和模型标识；未将模型生成的叙述当成执行日志，也不将该投影标称为官方网页原生导出。投影字节哈希已与 summary 复核一致。
\section{一条完整的核心工具链}
选择复杂仓库 \texttt{block\_route\_2}：
\begin{center}\begin{tabularx}{\linewidth}{p{18mm}p{37mm}X}\toprule
调用 / 结果 & 工具 & 核查内容\\\midrule
216 / 224 & \texttt{wg\_run\_case} & 在堵路场景执行原始策略，生成路径\\
239 / 243 & \texttt{wg\_audit\_path} & 原始路径独立验收不通过\\
256 / 260 & \texttt{wg\_recover\_case} & 按反馈刷新地图并重新规划\\
273 / 277 & \texttt{wg\_audit\_path} & 恢复路径独立验收通过\\\bottomrule
\end{tabularx}\end{center}
完整输入输出见\textbf{核心工具调用链.json}。每组调用与结果通过 \texttt{toolUseId} 匹配，来源为运行 A 的事件投影；内容未重新生成。
\section{Agnes 模型参与核心任务}
调用前的模型请求事件序号为 212、235、254、271。\par
\textbf{模型标识：}\texttt{agnes-3.0-flash}\par
\textbf{调用路由：}\texttt{account-agnes-cn-hackathon}\par
这些真实事件连接模型请求、工具调用与结果，展示模型参与测试与反馈恢复的过程。

调用方式为官方 AGH Host 的模型路由与受限 MCP 工具接口；模型编排动作，确定性规划器生成路径，几何验收器给出通过或失败。模型请求计数包含辅助请求；完整 token 与费用未知，原记录保留 null。
\section{关键配置与专业验证}
圆形机器人半径 0.18 m，安全余量 0.05 m；网格 0.25 m；验收器版本 \texttt{capsule-raw-grid-1}。同一目标与约束下，堵路原始路径最小剩余净空为 $-0.230$ m（FAIL），恢复后为 $0.145$ m（PASS）；路径长度为 41.5 m 与 48.0 m。

附件\textbf{关键配置.json}提供预算及运行设置；附件 05 提供三类样例、规范化 JSON 哈希及重放结果，三例均 \texttt{matches=true}。重放为离线确定性复核，模型请求为 0；保留原始失败，不以恢复通过覆盖初始问题。

\end{document}
